% Root proof with full independent actual-file audit by descent_groups.
\subsection{Every algebraic extension of a finite field}

Fix an odd prime $p$ and an algebraic closure $\overline{\mathbb F}_p$
with coherent finite-field arithmetic. Let
$\mathbb F_p\subseteq k\subseteq\overline{\mathbb F}_p$ be any
algebraic extension, finite or infinite. Define its supernatural degree
\[
 N_k=\prod_\ell\ell^{e_\ell},\qquad
 e_\ell=\sup\{v_\ell(n):\mathbb F_{p^n}\subseteq k\}
           \in\mathbb N\cup\{\infty\}.
\]
Divisibility by a finite integer and $\gcd(n,N_k)$ have their
prime-by-prime meanings. We put $h_2(k)=e_2$ and
\[
 T_k=\begin{cases}0,&h_2(k)=\infty,\\
                  \mathbb Z/2,&h_2(k)<\infty.
       \end{cases}                                         \tag{C.3.F1}
\]
The datum $h_2(k)$ is not assumed decidable merely from a program
enumerating the finite subfields.

\begin{lemma}[Proved: complete residual arithmetic over locally finite fields]
\label{C:algebraic_field_arithmetic}
The field $k$ consists exactly of the finite subfields
$\mathbb F_{p^n}$ with $n\mid N_k$. Every finite-dimensional
division algebra over $k$ is a field. For every finite extension
$K/k$, there is a canonical identification of groups
\[
                    K^\times/(K^\times)^2\simeq T_k.        \tag{C.3.F2}
\]
In the nonzero case the nontrivial element means the nonsquare class.
Nonsingular symmetric forms over $K$ are classified by rank and
this one bit. Alternating forms are classified by half their rank.
Hermitian forms for any nontrivial involution on such a $K$ are
classified by rank.
\end{lemma}
\begin{proof}
The degrees of finite subfields contained in $k$ are closed under
divisors and finite least common multiples. For a finite integer
$n\mid N_k$, choose a finite subfield realizing each of the finitely
many prime-power requirements of $n$; their compositum contains
$\mathbb F_{p^n}$. This proves the first assertion and the uniqueness
of the subfield specified by $N_k$.

Let $D$ be a finite-dimensional division algebra over $k$ and choose
a $k$-basis containing $1$. Its multiplication table uses only
finitely many coefficients, contained in a finite subfield $F$.
The $F$-span $D_0$ of this basis is a finite unital subring of $D$.
It has no zero divisors. Multiplication by a nonzero element is an
injective map of the finite set $D_0$ to itself, hence surjective;
so $D_0$ is a finite division ring. It is commutative by
\cite[Theorem III.4.4, p.29]{C-ArtinFiniteDivision}.
Its basis therefore commutes in $D$, so all of $D$ is commutative.

Any finite extension $K/k$ is a subfield of the same algebraic
closure and is a compositum $k\mathbb F_{p^n}$ for some finite $n$.
Its supernatural degree is $\operatorname{lcm}(N_k,n)$.
Consequently $h_2(K)$ is infinite exactly when $h_2(k)$ is infinite.
Over finite fields of odd order, a nonsquare becomes a square
precisely on an extension of even degree. If $h_2(K)=\infty$,
every element lies in a finite subfield admitting an even-degree
extension inside $K$, so every nonzero element is a square.
If $h_2(K)=a<\infty$, choose a finite subfield whose degree has
2-adic valuation $a$. A nonsquare in it remains nonsquare in all
larger finite subfields of $K$, since those extension degrees are
odd. Every element of $K$ can be put in such a larger subfield.
Finite-field square classes then show that there are exactly two
classes, proving (C.3.F2).

Every symmetric form and every finite list of its coefficients are
defined over a finite subfield. If two forms have the same rank
and determinant square class over $K$, also include a square root
of the determinant ratio in a finite subfield. The finite-field
classification then gives an isometry there. Necessity is the
determinant identity; thus the same classification holds over $K$.
The finite-field form theorem used here is the one checked in
Theorem~\ref{C:Om-finite-field-algorithm}. Alternating forms admit
a symplectic basis over any field of characteristic different from two.

For a nontrivial involution $\sigma$ of $K$, put $K_0=K^\sigma$.
Then $K/K_0$ is quadratic. To solve $N_{K/K_0}(x)=c$ for a given
$c\in K_0^\times$, choose a finite subfield $F_0\subset K_0$
containing $c$ and the coefficients of the minimal polynomial of
some $\alpha\in K\setminus K_0$. That polynomial is still
irreducible over $F_0$, and $F_0(\alpha)/F_0$ is quadratic.
The norm on finite fields is surjective, so it already gives
the required solution. Hermitian Gram--Schmidt and this norm
surjectivity classify hermitian forms by rank. A nonzero
$c$ with $\sigma(c)=-c$ similarly converts skew-hermitian forms
to hermitian forms. This proves all assertions.
\end{proof}

\begin{theorem}[Proved: the entire split orthogonal ring over every algebraic finite-field extension]
\label{C:all_algebraic_fields_WG}
Let $\mathscr I$ be the explicit absolute canonical index set over
$\overline{\mathbb F}_p$ of
Theorem~\ref{C:Om-canonical-templates}. For an index represented
by matrices over $\mathbb F_{p^n}$, its Galois orbit over $k$
is the finite orbit under coefficient Frobenius
\[
                    x\longmapsto x^{p^{\gcd(n,N_k)}},       \tag{C.3.F3}
\]
followed by canonicalization. These orbits, represented by their
least indices, give all indecomposable representations over $k$,
with Schur index one. The orbit is independent of the chosen
finite field of coefficients.

Use actual invariant forms and residue adjoints to partition the
self-dual indices into symmetric, alternating and unitary types
$\mathscr J_{\rm s},\mathscr J_{\rm a},\mathscr J_{\rm u}$,
and let $\mathscr J_{\rm h}$ choose one member of each other dual
pair, as in Proposition~\ref{C:C1_Galois_WG}. Then
\[
 WG_{\rm sp}(BO_m/k)\simeq
 \bigoplus_{i\in\mathscr J_{\rm s}}(\mathbb Z\oplus T_k)
 \ \oplus\!
 \bigoplus_{i\in\mathscr J_{\rm a}\cup\mathscr J_{\rm u}
                         \cup\mathscr J_{\rm h}}\mathbb Z .\tag{C.3.F4}
\]
The integer coordinates have the multiplicity conventions of
(C.3.Q2); the possible bit is the determinant square class of
the residual multiplicity form. This has no remaining division
algebra or hermitian-isometry classification input.

All ring products and exterior operations have terminating
procedures when the field is supplied as the indicated subfield
of the computable algebraic closure, together with decisions for
$n\mid N_k$ for finite $n$. Formula (C.3.F4) additionally uses
the explicit datum $h_2(k)\in\mathbb N\cup\{\infty\}$ to
decide whether its bit is present and to choose a uniform
nonsquare representative when it is present. No procedure for
recovering these data from an arbitrary computable presentation
of $k$ is asserted.
\end{theorem}
\begin{proof}
The finite-field intersection is
$k\cap\mathbb F_{p^n}=\mathbb F_{p^{\gcd(n,N_k)}}$.
The image of $\operatorname{Gal}(\overline{\mathbb F}_p/k)$
on $\mathbb F_{p^n}$ is therefore the cyclic subgroup generated
by (C.3.F3). The actual module, even when its canonical ordering
is not Galois invariant, is fixed by the $n$th Frobenius power.
Thus the finite orbit rule is valid and independent of $n$.
Theorem~\ref{C:perfect_Galois_indices} gives precisely one
indecomposable over $k$ for every orbit. Its residue division
algebra is a field by Lemma~\ref{C:algebraic_field_arithmetic},
so its Schur index is one. In particular its geometric dimension
is the size of the orbit times the dimension of any member.

Apply the checked hermitian-category decomposition to these
actual indecomposables. Their residue fields are finite over $k$.
Lemma~\ref{C:algebraic_field_arithmetic} gives exactly (C.3.F4),
including unitary and non-self-dual factors. It also proves that
the same bit convention works at every symmetric index, even
though their residue fields need not be identical. The underlying
representation multiplication is the orbit-sum rule of
Corollary~\ref{C:perfect_Green_image} with all indices equal to one.
The residual Gram determinants supply the additional possible bits;
geometric representation multiplicities alone do not supply them.

We verify the claimed effective access rather than assuming it.
The membership decisions enumerate all finite subfields of $k$;
finite-field names provide arithmetic, equality, and a surjective
enumeration of $k$. A finite input polynomial has its coefficients
in some $F=\mathbb F_{p^r}\subset k$. Factor it over $F$.
For a factor of degree $d$, all its roots lie in
$\mathbb F_{p^{rd}}$. On those roots the Galois image over $k$
is generated by the $\gcd(rd,N_k)$th Frobenius power.
Its finite orbits give the irreducible factors over $k$ by
multiplying the corresponding linear factors. Those coefficients
belong to $k$ and their finite-field names are explicit.
Thus factorization over $k$ terminates.

For a specified finite extension $K=k[z]/(f)$, choose a root of
$f$ in a finite field containing its coefficient field. Its
irreducibility over $k$ makes this a faithful embedding of $K$
into $\overline{\mathbb F}_p$. Write the resulting compositum
as $k\mathbb F_{p^n}$. Membership in its finite subfields is
decidable from $\operatorname{lcm}(N_k,n)$ by the same finite
integer divisibility tests. The preceding factorization procedure
therefore works over every such $K$ as well.

Here is a finite-stage procedure that does not need to decide
whether $k$ is finite. For an orbit representative over
$\mathbb F_{p^n}$ put $g=\gcd(n,N_k)$ and restrict scalars to
$\mathbb F_{p^g}$. Apply the finite-field algorithm to this
actual representation. Its indecomposable summands are all
isomorphic; choose the canonical one $M_0$. If the orbit has
size $a$, its residue endomorphism field has degree $a$ over
$\mathbb F_{p^g}$. Because $ga\mid n$ and
$\gcd(ga,N_k)=g$, that residue field remains a field on
extension to $k$. Thus $(M_0)_k$ is indecomposable and is the
desired $M_\Omega$. Its invariant reference form and residue
adjoint are obtained by the finite-field form algorithm.
Self-duality is already reflected over $\mathbb F_{p^g}$ by
Lemma~\ref{C:NoetherDeuring_rational}.

For a general finite input, put all matrices in a finite subfield
$F=\mathbb F_{p^r}\subset k$ and compute its endomorphism
algebra $E_F$. Its radical and semisimple quotient are computed
by finite enumeration; write the quotient as a product of
matrix algebras with centers $\mathbb F_{p^{rd_j}}$.
Let $g_j=\gcd(rd_j,N_k)$ and
$b=\operatorname{lcm}_j g_j$. The finite field
$F'=\mathbb F_{p^b}$ is contained in $k$.
Over $F'$ each center has already split into exactly the
factors that it has over $k$, and each resulting field factor
remains a field after scalar extension to $k$.
Indeed its defining Frobenius orbit is governed by $g_j$;
the equality $\gcd(b,rd_j)=g_j$ follows from $g_j\mid b\mid N_k$.
The radical commutes with these perfect field extensions,
since its tensor ideal is nilpotent with semisimple quotient.
Primitive-idempotent decomposition over the finite field $F'$
therefore gives summands that stay indecomposable over $k$.
Noether--Deuring also reflects every comparison with a chosen
representative to a common finite field of definition.

Compute the duals and reference forms there and reduce the
actual multiplicity Gram matrices to the residue fields, as
in the explicit extraction (C.3.Q4). Their ranks and determinant
square tests give all the coefficients in (C.3.F4).
Tensor and exterior matrices are formed before this extraction;
their entries again lie in a finite subfield. Thus every step
uses a finite field, finitely many integer membership decisions,
and the proved canonical procedures. There is no search for a
future stage whose existence has not been decided.

For clarity, the extra bit test is explicit. If $h_2(k)=\infty$,
it is absent for every finite residue field extension. Otherwise
a residue field $K/k$ of degree $a$ has
$h_2(K)=h_2(k)+v_2(a)$. A nonsquare in
$\mathbb F_{p^{2^{h_2(K)}}}\subset K$ remains nonsquare in $K$
and supplies a fixed representative. Given any nonzero residual
determinant, put it and this representative in a common finite
subfield of maximal 2-height; the ordinary finite-field square
test decides its bit. All remaining inputs and outputs are finite
matrices. Virtual operations follow coefficientwise from the
usual exterior-series inverse. This proves the complete operation
claim with its exact access assumptions.
\end{proof}

\begin{corollary}[Proved: the scalar ring over every locally finite odd field]
\label{C:algebraic_fields_scalar}
For the trivial group $O_0$ the answer is
\[
 WG_{\rm sp}(k)\simeq
 \begin{cases}
 \mathbb Z,&h_2(k)=\infty,\\
 \mathbb Z[t]/(2t,t^2),&h_2(k)<\infty.
 \end{cases}                                              \tag{C.3.F5}
\]
In the second case $t=[\langle a\rangle]-1$ for any nonsquare
$a\in k$. For $r\ge1$ its full virtual exterior structure is
\[
 \lambda^r(n+bt)=\binom nr+
       b\binom{n-1}{r-1}t,
       \qquad n\in\mathbb Z,\quad b\in\{0,1\}.            \tag{C.3.F6}
\]
In the first case $\lambda^r(n)=\binom nr$.
\end{corollary}
\begin{proof}
Rank and determinant from Lemma~\ref{C:algebraic_field_arithmetic}
give the additive groups. If $a$ is a nonsquare, the forms
$\langle a,a\rangle$ and $\langle1,1\rangle$ have equal rank
and determinant, hence are isometric. Thus $2t=0$.
Also $[\langle a\rangle]^2=1$, so $t^2=-2t=0$.
The determinant proves $t\ne0$. Finally
\[
 \lambda_z(n+bt)=(1+z)^n
       \left(\frac{1+(1+t)z}{1+z}\right)^b
       =(1+z)^n+b t z(1+z)^{n-1}.
\]
This identity is valid as a formal series for negative $n$
as well, and gives every coefficient in (C.3.F6).
The square-trivial case has only the rank coordinate.
\end{proof}

A concrete determinant square test already follows from the finite
subfield-membership decisions, by the proved factorization procedure.
The separate global 2-height datum specifies the displayed bit group
and a uniform nonsquare generator; it is not needed to test one
given determinant.

The extra presentation data above are substantive. The computable
fields in Theorem~\ref{C:locally-finite-equality-obstruction}
need not supply finite-subfield membership or square-root decisions
uniformly. The formulas here classify all the abstract fields in
this family; they do not contradict that decision obstruction.
